
\documentclass{article}
\usepackage{colortbl}
\usepackage{makecell}
\usepackage{multirow}
\usepackage{supertabular}

\begin{document}

\newcounter{utterance}

\centering \large Interaction Transcript for game `bbh', experiment `logical\_deduction\_five\_objects', episode 216 with refineonly{-}llama{-}3.1{-}8b{-}instruct.
\vspace{24pt}

{ \footnotesize  \setcounter{utterance}{1}
\setlength{\tabcolsep}{0pt}
\begin{supertabular}{c@{$\;$}|p{.15\linewidth}@{}p{.15\linewidth}p{.15\linewidth}p{.15\linewidth}p{.15\linewidth}p{.15\linewidth}}
    \# & \multicolumn{2}{c}{Player} && \multicolumn{2}{c}{Game Master} \\
    \hline

    \theutterance \stepcounter{utterance}  
    & & & \multicolumn{4}{p{0.6\linewidth}}{
        \cellcolor[rgb]{0.9,0.9,0.9}{
            \makecell[{{p{\linewidth}}}]{
                \texttt{\tiny{[P1$\langle$GM]}}
                \texttt{A logical deduction task which requires deducing the order of a sequence of objects.} \\
\\ 
\texttt{Q: The following paragraphs each describe a set of three objects arranged in a fixed order. The statements are logically consistent within each paragraph. In a golf tournament, there were three golfers: Amy, Eli, and Eve. Eve finished above Amy. Eli finished below Amy.} \\
\texttt{Options:} \\
\texttt{(A) Amy finished last} \\
\texttt{(B) Eli finished last} \\
\texttt{(C) Eve finished last} \\
\texttt{A: (B)} \\
\\ 
\texttt{Q: The following paragraphs each describe a set of three objects arranged in a fixed order. The statements are logically consistent within each paragraph. On a shelf, there are three books: a white book, a green book, and an orange book. The green book is to the right of the white book. The orange book is the rightmost.} \\
\texttt{Options:} \\
\texttt{(A) The white book is the leftmost} \\
\texttt{(B) The green book is the leftmost} \\
\texttt{(C) The orange book is the leftmost} \\
\texttt{A: (A)} \\
\\ 
\texttt{Q: The following paragraphs each describe a set of three objects arranged in a fixed order. The statements are logically consistent within each paragraph. On a shelf, there are three books: a red book, a gray book, and a white book. The white book is to the left of the gray book. The red book is the second from the left.} \\
\texttt{Options:} \\
\texttt{(A) The red book is the leftmost} \\
\texttt{(B) The gray book is the leftmost} \\
\texttt{(C) The white book is the leftmost} \\
\texttt{A: (C)} \\
\\ 
\texttt{Q: The following paragraphs each describe a set of five objects arranged in a fixed order. The statements are logically consistent within each paragraph. In a golf tournament, there were five golfers: Eve, Amy, Ada, Rob, and Joe. Amy finished second{-}to{-}last. Rob finished below Eve. Ada finished above Joe. Joe finished second.} \\
\texttt{Options:} \\
\texttt{(A) Eve finished second} \\
\texttt{(B) Amy finished second} \\
\texttt{(C) Ada finished second} \\
\texttt{(D) Rob finished second} \\
\texttt{(E) Joe finished second} \\
\texttt{A:} \\
            }
        }
    }
     \\ \\

    \theutterance \stepcounter{utterance}  
    & & \multicolumn{4}{p{0.6\linewidth}}{
        \cellcolor[rgb]{0.95,0.95,0.95}{
            \makecell[{{p{\linewidth}}}]{
                \texttt{\tiny{[P1$\leftrightarrow$P]}}
                \texttt{{[}DRAFT{]}} \\
\texttt{To deduce the correct order, let's analyze the given information:} \\
\\ 
\texttt{1. Amy finished second{-}to{-}last.} \\
\texttt{2. Rob finished below Eve.} \\
\texttt{3. Ada finished above Joe.} \\
\texttt{4. Joe finished second.} \\
\\ 
\texttt{From statement 4, we know Joe finished second. Since Amy finished second{-}to{-}last (statement 1), Amy must be in the fourth position. } \\
\\ 
\texttt{Now, let's consider the remaining positions. Since Joe is in the second position, the order so far is: Joe, Amy (second{-}to{-}last).} \\
\\ 
\texttt{From statement 3, we know Ada finished above Joe. This means Ada must be in the first position, as Joe is in the second position.} \\
\\ 
\texttt{Now, let's consider the remaining positions. We have Ada in the first position, Joe in the second position, and Amy in the fourth position. The only remaining position is the third position, which must be occupied by Eve or Rob.} \\
\\ 
\texttt{From statement 2, we know Rob finished below Eve. This means Eve must be in the third position, as Rob is below Eve.} \\
\\ 
\texttt{Now, let's summarize the order: Ada, Joe, Eve, Amy, Rob.} \\
\\ 
\texttt{The correct answer is not among the options provided. However, based on the analysis, the correct order is: } \\
\\ 
\texttt{Ada, Joe, Eve, Amy, Rob} \\
            }
        }
    }
    & \\ \\

    \theutterance \stepcounter{utterance}  
    & \multicolumn{4}{p{0.6\linewidth}}{
        \cellcolor[rgb]{0.561,0.737,0.561}{
            \makecell[{{p{\linewidth}}}]{
                \texttt{\tiny{[P1$\rangle$GM]}}
                \texttt{The following paragraphs each describe a set of five objects arranged in a fixed order. The statements are logically consistent within each paragraph. In a golf tournament, there were five golfers: Eve, Amy, Ada, Rob, and Joe. Amy finished second{-}to{-}last. Rob finished below Eve. Ada finished above Joe. Joe finished second.} \\
\\ 
\texttt{Options:} \\
\texttt{(A) Eve finished second} \\
\texttt{(B) Amy finished second} \\
\texttt{(C) Ada finished second} \\
\texttt{(D) Rob finished second} \\
\texttt{(E) Joe finished second} \\
            }
        }
    }
    & & \\ \\

    \theutterance \stepcounter{utterance}  
    & & & \multicolumn{2}{p{0.3\linewidth}}{
        \cellcolor[rgb]{0.9,0.9,0.9}{
            \makecell[{{p{\linewidth}}}]{
                \texttt{\tiny{[GM$|$GM]}}
                \texttt{Target: (E)} \\
            }
        }
    }
    & & \\ \\

    \theutterance \stepcounter{utterance}  
    & & & \multicolumn{2}{p{0.3\linewidth}}{
        \cellcolor[rgb]{0.9,0.9,0.9}{
            \makecell[{{p{\linewidth}}}]{
                \texttt{\tiny{[GM$|$GM]}}
                \texttt{game\_result = LOSE} \\
            }
        }
    }
    & & \\ \\

\end{supertabular}
}

\end{document}
